% =====================================================================
\chapter{Proposiciones, conectivas y tablas de verdad}
\label{tema:1}
% =====================================================================

Una demostración es una cadena de afirmaciones enlazadas por reglas de razonamiento. Antes de estudiar las técnicas de demostración conviene fijar qué es una afirmación, cómo se combinan varias mediante conectivas y cuándo dos combinaciones dicen lo mismo. Este tema desarrolla la lógica proposicional con el enfoque del capítulo~2 de~\cite{hammack} (§2.1–2.6 y §2.11) y del capítulo~1 de~\cite{velleman}; la parte algebraica sigue la §3.4 de~\cite{mcs}.

\section{Proposiciones y conectivas}

\begin{definicion}[Proposición]\label{def:proposicion}
Una \emph{proposición} es un enunciado del que tiene sentido afirmar que es verdadero o que es falso, y que es exactamente una de las dos cosas. Su \emph{valor de verdad} es $\V$ (verdadero) o $\F$ (falso).
\end{definicion}

\addtocounter{definicion}{1}%

A partir de proposiciones $P$ y $Q$ se forman proposiciones compuestas mediante \emph{conectivas lógicas}. El valor de verdad de la proposición compuesta depende solo de los de $P$ y $Q$, y se describe con una \emph{tabla de verdad}.

\begin{definicion}[Negación, conjunción y disyunción]\label{def:conectivas}
Sean $P$ y $Q$ proposiciones.
\begin{enumerate}
  \item La \emph{negación} $\neg P$ («no $P$») es verdadera si y solo si $P$ es falsa.
  \item La \emph{conjunción} $P\wedge Q$ («$P$ y $Q$») es verdadera si y solo si $P$ y $Q$ son ambas verdaderas.
  \item La \emph{disyunción} $P\vee Q$ («$P$ o $Q$») es verdadera si y solo si al menos una de las dos es verdadera.
\end{enumerate}
\end{definicion}

\begin{table}[h]
\centering
\caption{Tablas de verdad de las conectivas básicas.}\label{tab:conectivas}
\begin{tabular}{@{}cc|ccccc@{}}
\toprule
$P$ & $Q$ & $\neg P$ & $P\wedge Q$ & $P\vee Q$ & $P\Rightarrow Q$ & $P\Leftrightarrow Q$ \\
\midrule
$\V$ & $\V$ & $\F$ & $\V$ & $\V$ & $\V$ & $\V$ \\
$\V$ & $\F$ & $\F$ & $\F$ & $\V$ & $\F$ & $\F$ \\
$\F$ & $\V$ & $\V$ & $\F$ & $\V$ & $\V$ & $\F$ \\
$\F$ & $\F$ & $\V$ & $\F$ & $\F$ & $\V$ & $\V$ \\
\bottomrule
\end{tabular}
\end{table}

La disyunción matemática es \emph{inclusiva}: $P\vee Q$ es verdadera también cuando lo son las dos. La disyunción \emph{exclusiva}, «o $P$ o $Q$, pero no ambas», se escribe $P\xor Q$ y se estudia en el \cref{eje:1-xor}.

\section{Condicional y bicondicional}

\begin{definicion}[Condicional]\label{def:condicional}
El \emph{condicional} $P\Rightarrow Q$ («si $P$, entonces $Q$») es falso únicamente cuando $P$ es verdadera y $Q$ es falsa. $P$ se llama \emph{hipótesis} o antecedente, y $Q$, \emph{tesis} o consecuente.
\end{definicion}

La tabla del condicional (\cref{tab:conectivas}) se entiende como una promesa: «si apruebas el examen, te invito a cenar» solo se incumple si se aprueba y no hay cena. Si no se aprueba, la promesa no se ha roto, haya cena o no. Por eso un condicional con hipótesis falsa es verdadero; se dice que es \emph{vacuamente} verdadero.

El condicional $P\Rightarrow Q$ admite varias lecturas equivalentes, que conviene reconocer porque aparecen constantemente en los enunciados de los teoremas~\cite[§2.3]{hammack}:
\begin{center}
\begin{tabular}{@{}l@{}}
«si $P$, entonces $Q$»; \quad «$Q$ si $P$»; \quad «$P$ solo si $Q$»; \\
«$P$ es condición suficiente para $Q$»; \quad «$Q$ es condición necesaria para $P$».
\end{tabular}
\end{center}

\begin{definicion}[Bicondicional]\label{def:bicondicional}
El \emph{bicondicional} $P\Leftrightarrow Q$ («$P$ si y solo si $Q$») es verdadero cuando $P$ y $Q$ tienen el mismo valor de verdad. Se lee también «$P$ es condición necesaria y suficiente para $Q$».
\end{definicion}

\begin{definicion}[Recíproco, contrarrecíproco e inverso]\label{def:reciproco}
Dado el condicional $P\Rightarrow Q$, se llaman \emph{recíproco} a $Q\Rightarrow P$, \emph{contrarrecíproco} a $\neg Q\Rightarrow\neg P$ e \emph{inverso} a $\neg P\Rightarrow\neg Q$.
\end{definicion}

\addtocounter{definicion}{1}%

\section{Fórmulas, tautologías y contradicciones}

\begin{definicion}[Fórmula proposicional]
Las \emph{fórmulas proposicionales} sobre las variables $P_1,\dots,P_n$ son las expresiones que se obtienen a partir de las variables aplicando un número finito de veces las conectivas $\neg$, $\wedge$, $\vee$, $\Rightarrow$ y $\Leftrightarrow$, con paréntesis para indicar el orden. Una \emph{asignación} da a cada variable un valor de verdad, y la tabla de verdad de la fórmula lista su valor para cada asignación.
\end{definicion}

\begin{proposicion}\label{prop:filas}
Una fórmula con $n$ variables tiene una tabla de verdad de $2^{n}$ filas.
\end{proposicion}

\begin{definicion}[Tautología y contradicción]
Una fórmula es una \emph{tautología} si es verdadera para todas las asignaciones, y una \emph{contradicción} si es falsa para todas.
\end{definicion}

\addtocounter{definicion}{1}%

El razonamiento del ejemplo anterior evita escribir la tabla completa: para comprobar que un condicional es tautología basta ver que la fila «hipótesis verdadera y tesis falsa» es imposible. Esta forma de pensar es ya una demostración por reducción al absurdo (\cref{tema:4}).

\section{Equivalencia lógica}

\begin{definicion}[Equivalencia lógica]\label{def:equivalencia}
Dos fórmulas $\Phi$ y $\Psi$ en las mismas variables son \emph{lógicamente equivalentes}, y se escribe $\Phi\eqlog\Psi$, si tienen el mismo valor de verdad en todas las asignaciones; equivalentemente, si $\Phi\Leftrightarrow\Psi$ es una tautología.
\end{definicion}

\begin{proposicion}[Leyes lógicas]\label{prop:leyes}
Para proposiciones cualesquiera $P$, $Q$ y $R$:
\begin{enumerate}
  \item doble negación: $\neg\neg P\eqlog P$;
  \item leyes de De Morgan: $\neg(P\wedge Q)\eqlog\neg P\vee\neg Q$ \ y \ $\neg(P\vee Q)\eqlog\neg P\wedge\neg Q$;
  \item distributivas: $P\wedge(Q\vee R)\eqlog(P\wedge Q)\vee(P\wedge R)$ \ y \ $P\vee(Q\wedge R)\eqlog(P\vee Q)\wedge(P\vee R)$;
  \item condicional como disyunción: $P\Rightarrow Q\eqlog\neg P\vee Q$;
  \item contrarrecíproco: $P\Rightarrow Q\eqlog\neg Q\Rightarrow\neg P$;
  \item negación del condicional: $\neg(P\Rightarrow Q)\eqlog P\wedge\neg Q$;
  \item bicondicional: $P\Leftrightarrow Q\eqlog(P\Rightarrow Q)\wedge(Q\Rightarrow P)$.
\end{enumerate}
Además, $\wedge$ y $\vee$ son conmutativas y asociativas.
\end{proposicion}

Las leyes de la \cref{prop:leyes} permiten transformar fórmulas sin construir tablas, igual que las propiedades de la suma y el producto permiten manipular expresiones algebraicas~\cite[§3.4]{mcs}.

\addtocounter{definicion}{1}%

\section{Razonamientos válidos}\label{sec:inferencia}

\begin{definicion}[Razonamiento válido]
Un razonamiento con premisas $\Phi_1,\dots,\Phi_k$ y conclusión $\Psi$ es \emph{válido} si $(\Phi_1\wedge\dots\wedge\Phi_k)\Rightarrow\Psi$ es una tautología, es decir, si en toda asignación que hace verdaderas todas las premisas también es verdadera la conclusión.
\end{definicion}

\begin{proposicion}[Reglas de inferencia]\label{prop:inferencia}
Los siguientes razonamientos son válidos~\cite[§2.11]{hammack}:
\begin{center}
\begin{tabular}{@{}lll@{}}
\toprule
Nombre & Premisas & Conclusión \\
\midrule
\emph{Modus ponens} & $P\Rightarrow Q$, \ $P$ & $Q$ \\
\emph{Modus tollens} & $P\Rightarrow Q$, \ $\neg Q$ & $\neg P$ \\
Silogismo hipotético & $P\Rightarrow Q$, \ $Q\Rightarrow R$ & $P\Rightarrow R$ \\
Silogismo disyuntivo & $P\vee Q$, \ $\neg P$ & $Q$ \\
\bottomrule
\end{tabular}
\end{center}
\end{proposicion}

\addtocounter{definicion}{1}%

\addtocounter{definicion}{1}%

